% Bare circuit copied from the paper's fig:distill-overview-bare-intro.
\begin{frame}{Distillation: The 4-to-1 Circuit}
  \centering
  \begin{tikzpicture}
    \path[use as bounding box] (-6.6,-3.15) rectangle (6.6,3.05);
    \node[inner sep=0pt,text=CVQCInk] at (0,0.05) {%
      \scalebox{1.5}{%
        % qcircuit's meterB uses a text-size command inside math mode.
        \let\tiny\scriptstyle
        \hspace*{0.55cm}% Include the input labels in the visual bounds.
        $\Qcircuit @C=1.4em @R=1.2em {
          \lstick{\ket{0}} & \qw & \ctrl{1} & \targ & \qw & \qw \\
          \lstick{\ket{0}} & \qw & \targ & \qw & \meter & \cw \\
          \lstick{\ket{0}} & \qw & \ctrl{1} & \ctrl{-2} & \meterB{X} & \cw \\
          \lstick{\ket{0}} & \qw & \targ & \qw & \meter & \cw
        }$%
      }%
    };
  \end{tikzpicture}

  \note{
    \footnotesize
    This is the paper's basic 4-to-1 distillation circuit. The first wire
    survives; wires two, three, and four are measured in the Z, X, and Z
    bases, respectively. C-dist denotes the three-CNOT unitary part,
    excluding these measurements: CNOT from one to two, from three to
    four, and from three to one.
    Begin with the plaintext ancilla example. On four zero inputs the
    CNOTs preserve the inputs. Check that the two Z-basis outcomes are
    zero; either X-basis outcome is accepted, since that measurement is
    random on a zero state. Honest execution then accepts and leaves zero
    on the first wire. With at most one bad input, the routine either
    rejects or leaves the intended surviving zero state. This statement
    concerns honest execution of the circuit. Handling deviations during
    authenticated execution is the next task. The manuscript also extends
    the pattern to EPR pairs and Hadamard magic states.
    Source: Technical Overview, basic 4-to-1 circuit; Distillation section,
    accepted measurement results and correctness analysis.
  }
\end{frame}
